\chapter{Goodness of Fit}

Congrats, you've made it this far, and not much more separates you from the back cover of this book\sidenote{Admittedly, there are quite a few pages, but the real "meat" of the course is behind us}. We wrap up our general review of inferences with this final type that differs from all the others in an important way: we want to test if the \emph{pdf} we assume the data to be extracted from is actually the "correct one". Well, without getting ahead of ourselves we want to quantify, in some statistically correct way, \textbf{how much the observed data is in accord with the assumed distribution}. 

Up until now we have assumed that the observed data come from a known \emph{pdf} (or class of \emph{pdf's}) depending on some parameter, and we have worked long and hard to try to say \emph{something} about that parameter. Well, now we must zoom out a step, and justify why we are assuming that distribution in the first place. Without this initial justification, whatever we end up discovering about the parameters that govern $p(x)$ must be thrown directly in the bin, as if we have failed to convince ourselves that $x$ is distributed following $p$ in the first place then all our work is for naught.

Naturally, this being the last chapter does not mean that it is the last thing we wish to verify during our analysis. Quite the contrary, it can be wise to first verify how well our data is in accord with the assumed \emph{pdf}, and to then analyse it, as to not waste precious resources and time.

Convinced? Let's get into it:

GOF tests are in many ways similar to hypothesis tests, and once again we choose a Test statistic $T(x)$ and find a threshold $T(x)>q_\alpha$ such that the Test has a significance of $\alpha$ (chosen \emph{a priori}).

In this case, however, only the null hypothesis is specified (simple or composite as it may be), and the alternative is simply "anything else". By anything we really do mean it, as an alternative could be that the data does not even follow a \emph{pdf} (the phenomenon is \emph{non-probabilistic}). Therefore, the alternative space is not defined, and we can't define a power function, much less talk about UMP Tests.

Having said this, we can still talk about Tests that are more "powerful" (in quotes) than others, as this non-definition will include subjective decisions on which Test to prefer.

We struggle even to define the concept of \emph{consistency} in this case, and typically fall upon guaranteeing that the GOF Test be at least \emph{unbiased} on the alternatives that we are able to think of.

Not having defined alternatives, we can't talk about \emph{priors} (there may not even be a \emph{pdf} that follows the data), which is a problem for the Bayesian school, meaning that the very concept of GOF Tests are, at their core, frequentist. To round things out, it also strongly violates the Likelihood principle, as the Likelihood gives no useful information in this formulation.

\section{p-values}

The statistic at the core of almost every GOF Test is the famous (or infamous, if you were not a fan of the Lab courses) \textbf{p-value} statistic. This is perhaps one of the most common and commonly misunderstood concepts in data analysis, and so we take our time in defining it. Let's begin with an example:

Imagine flipping a coin 10 times and getting 10 heads in a row. You suspect the coin is unfair, and thus construct a null hypothesis to hopefully reject $H_0$: The coin is fair (the outcomes of $n$ flips follows a Binomial with $p=0.5$). Our alternative hypothesis is that the coin is \emph{unfair}, as in, anything that is not $H_0$, and thus a goodness-of-fit test must be used. If the coin really were fair, seeing 10 heads would be extremely rare. The $p$-value quantifies exactly how "rare" or "surprising" the observed data is, assuming $H_0$. We can see the parallels with the previous chapter, and we expect to use $p$ as a sort of Test statistic with a certain critical region. Therefore, it is intuitive that $p$ depend on the data $x$, and in some way must be combined with $H_0$. 

\begin{definition}[p-value]
	A \emph{p-value} is a statistic $p(x)$ such that $p\sim U[0,1]$ under $H_0$
\end{definition}

Mathematically, if $T(x)$ is our Test statistic and $F_T$ is its cumulative distribution function under $H_0$, then $p = 1 - F_T(t_{obs})$\sidenote{This assumes a right (upper) tailed test, where large values of $T$ suggest evidence against $H_0$. For left (lower) tailed tests, $p = F_T(t_{obs})$}, where $t_{obs} = T(x_{obs})$, the sampled Test statistic. Because the \emph{CDF} of any continuous random variable is uniformly distributed, as the probability integral transform theorem \ref{th:universal_uniform} states, the p-value follows $U[0, 1]$.

Essentially, it is a function of the data such that its distribution is uniform between $0$ and $1$ and such that it is small for hypotheses distant from $H_0$.

The goodness of fit test, in this case, corresponds to the critical region $p<\alpha$, otherwise stated as $P(p<\alpha)=\alpha\;\forall\alpha$.

We also wish for the GOF test to give unbiased results, and so we need for this statistic $p(x)$ to move towards lower values in every hypothesis in which $H_0$ is not verified (for right-tailed tests, as the opposite can be true granted we use a left-tailed test). Synthetically:

	\[
	F_{\cancel{H_0}}(p) \geq F_{H_0}(p)
	\]

This gives the $p$-value the interpretation of "probability" (under $H_0$) that the experiment give a result that deviates further from $H_0$ that what was observed. 

\begin{figure}[H]
	\centering
	\input{images/p_value_plot.pgf}
	\caption{Distribution of the $p$-value statistic under the null hypothesis $H_0$ and an example alternative $H_1$. A vertical line corresponding to $\alpha=0.05$ significance level is also shown. The probability of a uniformly distributed value to fall below the threshold, which corresponds to the probability of committing a Type I Error, is $\alpha$.}
	\label{images/p_value_plot}
\end{figure}

From figure \ref{images/p_value_plot} we can see the utility in our construction: all values of $p$ are equally likely under $H_0$, while for $H_1$ lower values are more probable. Thinking about it in a frequentist sense, if we were to repeat many trials and extract a $p$-value for each of them, under $H_0$ no specific value holds more importance than any other, and we expect to get all of them sooner or later. Under $H_1$, lower values are more likely. Therefore, our test corresponds to the critical region $p<\alpha$, and only in that case can we reject the null hypothesis. 

We can also visualize the different types of errors we may commit: obtaining $p<\alpha$ suggests that it was not extracted from $H_0$, though naturally there is $\alpha$ chance for it to come from the uniform distribution, and so a probability of $\alpha$ to have committed a Type I Error, and erroneously rejected the null hypothesis. Hopefully this example clears the question of why we can attribute to the $p$-value the probability under $H_0$ to have a more extreme solution than what was observed.

\subsubsection{What the p-value is not:}
\begin{itemize}
	\item It is \textbf{NOT} the probability that the null hypothesis is true. ($P(H_0 | \text{data}) \neq \text{p-value}$)
	\item It is \textbf{NOT} the probability that the results were caused by chance
	\item It is \textbf{NOT} the significance of the test, as that is the name given to $\alpha$, which has been priorly fixed.  
\end{itemize}

Let's also take this opportunity, while we are talking about incorrect interpretations, to introduce some cases for which the abuse of goodness-of-fit reasoning brings erroneous conclusions. First of all, let's keep in mind that the alternative hypothesis in this case is "anything that is not $H_0$". Therefore, to make a conceptual analogy, if our null hypothesis is "God does not exist" and we reject it using a GOF Test, it does not directly follow that God does indeed exist. It means that the specific data we collected is incompatible with a model where no such influence exists. In a goodness-of-fit context, rejecting $H_0$ only tells you that your null model is a poor description of reality; \emph{it doesn’t automatically validate a specific alternative}. There could be a third, unconsidered variable, a flaw in the experimental design, or simply a Type I Error. 

\begin{example}
	Perhaps a less philosophical example: suppose that we want to test the fairness of a six-sided die. Our null hypothesis is that the die rolls each number $1-6$ with an equal $\frac{1}{6}$ probability. We decide to test this, roll $100$ times, and only get even numbers for all of them. Can we say that the die only rolls even numbers? No, we cannot, all we can do, with this GOF formalism, is \emph{reject the hypothesis that the die is fair}, though we cannot justify our theory as to \emph{why} it is not fair.
\end{example}

From these examples we can conclude one thing: \textbf{any "discoveries" based only on GOF tests cannot be accepted}. The goodness-of-fit is a useful and necessary tool that must be implemented to justify any \emph{pdf} we may be inclined to assume, but it absolutely cannot be used as a stand-alone decision-making tool. 

\subsubsection{Combination of p-values}

Let's say that in two iterations of the same experiment I measure two different $p$-values $p_1$ and $p_2$. We'd like a way to express the combination of these two to have one total $p$-values indicative of both experiments. 

Since we have attributed the definition of probability to $p$-values, and the two experiments are considered to be independent, can we simply take the product of these two $p$-values and combine them as we would probabilities of i.i.d random variables? The idea is to see if $P\left(x\leq\alpha\right) = P\left(p_1\leq\alpha\right)\cdot P\left(p_2\leq\alpha\right)$ where $x=p_1\cdot p_2$.

Well, we want $x$\sidenote{Please allow us to introduce a new letter for the sake of not calling $p=p_1\cdot p_2$, and using the letter "p" once again, even though $x$ seems out of place} to be distributed as a $p$-value should, which is following a standard uniform distribution. To see this, we wish to calculate the cumulative of the product $x$:

	\[
	F_X(x) =  P(p_1 \cdot p_2 \leq x)
	\]
To find this, we integrate the joint PDF (which is $1 \cdot 1 = 1$) over the region in the unit square where $p_2 \leq \frac{x}{p_1}$:


\begin{align*}
	F_X(x) &= \int_{0}^{x} 1 \, \dd p_1 + \int_{x}^{1} \frac{x}{p_1} \dd p_1 = x + x \left[ \ln\left(p_1\right) \right]_{x}^{1}\\
	&= x - x\ln(x) = x(1 - \ln(x))
\end{align*}


\begin{figure}[H]
	\centering
	\input{images/p_value_comb.pgf}
	\caption{Integration bounds for two variables $p_1$, $p_2$ uniformly distributed between $0$ and $1$, with the line $p_2=\frac{x}{p_1}$ explicitly drawn. This line shows, for each $p_1$, what value $p_2$ must adopt for the product to be $\alpha=0.05$}
	\label{images/p_value_comb}
\end{figure}

As we can see, the cumulative of the joint \emph{pdf} is not simply $x$, and therefore does not describe a uniform distribution, and so simply multiplying the two $p$-values is not a valid way to combine them. It's interesting to see how this mistake can cost us our theory developed thus far, if I want to combine two $p$-values of $0.1$, I would incorrectly state that my combined $p$-value be equal to $0.01$ when I do the banal product. 

In reality, under the null hypothesis, the probability of getting a product $\leq 0.01$ is $0.01(1 - \ln(0.01)) \approx 0.056$, which means that I would claim a significant result at the $1\%$ level more than $5\%$ of the time! Clearly this is not the strategy.

What we have shown is that to combine two $p$-values, instead of multiplying them we combine them using the formula we derived, as the probability integral transform guarantees that the CDF of a variable transformation is distributed uniformly:

	\[
	p = p_1p_2 (1 - \ln(p_1p_2))
	\tag{Combination of two $p$-values}
	\]
	

This is not the only way that exists, and for $n$ $p$-values we prefer to calculate:

	\[
	q=-2\sum_i\ln\left(p_i\right) = -2\ln\left(\prod_i p_i\right)
	\]

To find the threshold $q>q_\alpha$, we want to find how $q_i=-2\ln\left(p_i\right)$ is distributed. Trough a variable change\sidenote{$q_i=-2\ln\left(p_i\right) \implies p_i=e^{-\frac{q_i}{2}}$}:
\begin{align*}
	\mathfrak{p}(q_i)&=\mathfrak{p}(p_i(q_i))\left|\frac{\dd q_i}{\dd p_i}\right|^{-1}_{p_i=p_i(q_i)}\\
	&=\frac{1}{2}e^{-\frac{q_i}{2}}
\end{align*}
Now we consider that:
	\[
	\chi^2_\nu = \frac{1}{2^{\frac{\nu}{2}}\Gamma\left(\frac{\nu}{2}\right)}x^{\frac{\nu}{2}}e^{-\frac{x}{2}} \Rightarrow \chi_2^2 \sim \frac{1}{2}e^{-\frac{1}{2}} \Rightarrow q_i=-2\ln\left(p_i\right) \sim \chi_2^2
	\]
Since $q=\sum_iq_i$, using the fact that the sum of independent chi-squared variables is again chi-squared, we get:
\[
	q\sim\chi^2_{2n}
\]
Therefore we can find the threshold $q>q_\alpha$ through the use of tables.

This combination of $p$-values is not just a didactic formalism; to be completely transparent, it is necessary to consider the ensemble of results from all repetitions of an experiment. For example, if I run an experiment an arbitrarily large number of times, sooner or later I will obtain at least one significant $p$-value purely by chance, due to the inherent probability of committing a Type I Error (which is exactly what $\alpha$ quantifies). 

It would be intellectually dishonest to report only that singular result\sidenote{This is a form of \emph{publication bias} or \emph{p-hacking}, where non-significant results are ignored, often leading to the "Look Elsewhere Effect" or the "File Drawer Problem"}, ignoring the others; in practice, this is no better than simply fabricating data. Therefore, two options come to mind: re-scale $\alpha$ for each trial (the \emph{Bonferroni Correction} used in life sciences), as to compensate for the "multiplicity" of many independent trials, or to combine all of the obtained $p$-values into one singular one. While the first option can be overly conservative, the second (combining them) follows naturally from the methods we have developed in this chapter and is the one most commonly used in our field.

\section{GOF Tests from the LR statistic}
With a solid understanding of $p$-values under our belt, it's time to move on to the practical applications of this theory. Let's begin with the case where the most general hypothesized \emph{pdf} can be described as a member of a parametrized space $p(x;\theta)$. This space can even have a high number of dimensions, and $H_0$ is the case of $\theta=\theta_0$. As we have done in the previous chapter, $H_0$ can have an arbitrary number of nuisance parameters $\nu$.

Following the formalism developed in \ref{subsec:LR_test}, we once again use the Log Likelihood Ratio (LLR):

	\[
	\lambda(x) = 2\ln \frac{\sup_{\theta,\nu}p(x;\theta,\nu)}{\sup_{\nu}p(x;\theta_0,\nu)}
	\]

Which, as we have done before, can be re-written by replacing the sup's with \emph{MLE}\sidenote{We have gotten used to this switch from sup to \emph{MLE}, but we still explicit the original supremum for edge cases where the \emph{MLE} does not exist (we can consider it to be an asymptotic \emph{MLE})}.

Once again, we know through Wilks' theorem that under $H_0$, $\lambda(x)$ will asymptotically follow a central $\chi^2_d$ distribution, where the number of degrees of freedom $d$ is equal to the number of parameters fixed by the null hypothesis (the dimension of $\theta$). For all cases where $\theta \neq \theta_0$, the distribution shifts towards higher values\todo{Why???} (following a non-central $\chi^2$), ensuring the test is unbiased.

Therefore, the Test with critical region $\lambda(x)>q_\alpha$ will be unbiased and right-tailed, so a good $p$-value is $p=1-F_{\chi^2_d}(\lambda)$. 

\subsection{Goodness of Fit applied to histograms}
In this section we will follow the formalism introduced in section \ref{sec:fitting_hist}, revisiting what we have discussed using the new tools we have at our disposal. Specifically, we will look at variable-$N$ histograms, as we cover in section \ref{subsec:variable_N_hist}.
Considering a sequence of experimental measurements of a real observable $x$, let $k_i$ be the number of events $\{x_j\}$ falling into the bin $i\in \left[1,\dots,n\right]$. Therefore, what we have is a histogram of $n$ bins, $N$ measurements, completely defined by $\{k_i\}_n$ such that $\sum_i^n k_i =N$. We want to estimate the parameters $\theta$ of the \emph{pdf} from which $x$ was extracted, which we know to be $p(x;\theta)$.

As we have already discussed, each bin count is independent and follows a Poisson distribution, such that:

	\[
	p(x;\theta) = \prod_i^n \mathcal{P}\left(k_i,\mu_i(\theta)\right)
	\]

What we are looking for, then, is a point estimation of $\theta$, and we'll use the \emph{MLE}:

	\[
	\hat{\theta}_{\text{MLE}} = \arg \max \prod_i^n \mathcal{P}\left(k_i;\mu_i(\theta)\right)
	\]

We know $\mu_i(\theta)=\E\left[k_i\right]$ to be the expected number of events in the $i$-th bin.

Well, now we are more mature than we were all those chapters ago, and instead of blindly accepting the $\hat{\theta}_{\text{MLE}}$ we would like to confirm that the observed distribution does not statistically differ from the hypothesized \emph{pdf}. We assume that if we add \emph{enough} parameters to our current \emph{pdf}, then we can generally obtain whichever vector of expected values $\mu_i(\theta)$. In a certain sense, we can't explicit the most general alternative hypothesis, but proceeding in this way we can convince ourselves that we could arbitrarily "fit any histogram" given enough parameters.

In this case, the LLR becomes:

	\[
	\lambda(x) = 2\ln\frac{\sup_{\theta,\nu}\prod_i\mathcal{P}\left(k_i;\mu_i(\theta,\nu)\right)}{\sup_\theta\prod_i\mathcal{P}\left(k_i;\mu_i(\theta)\right)}
	\]

To evaluate this, let's look at the numerator. The term $\sup_{\theta,\nu}\mu_i(\theta,\nu)$ represents the mean of the Poisson distribution in the $i$-th bin under the most general possible hypothesis. If we assume this general model has enough parameters to "fit any histogram" to any degree of precision, it becomes what is known as the \emph{Saturated Model}. 

In this saturated case, the Likelihood is maximized when each bin's expected value is exactly equal to its observed count, $\hat{\mu}_{i} = k_i$. Essentially, we are allowing each of the $n$ bins to behave as an independent free parameter. Conversely, the denominator represents our specific physical model, which constrains those $n$ bins using only a small set of parameters $\theta$.

Following Wilks' theorem, the degrees of freedom for our test statistic will be the difference in the number of free parameters between these two models. Since the Saturated Model uses $n$ parameters and our model uses $\dim(\theta)$, we find:
	\[
	DOF = n - \dim(\theta)
	\]

By evaluating the Poisson product at these respective maxima, which are $\hat{\mu}_{i} = k_i$ for the numerator and $\hat{\mu}_i(\hat{\theta})$ for the denominator, our expression becomes:\sidenote{Consider the "hats" to be the \emph{MLE}'s, not explicitly stated for the sake of clarity}

	\[
	\lambda(x) = 2\ln\frac{\prod_i^n \mathcal{P}\left(k_i;k_i\right)}{\prod_i^n \mathcal{P}\left(k_i;\hat{\mu}_i(\hat{\theta})\right)}
	\]
	
We expand using the definition of a Poisson statistic:

\begin{align*}
	\lambda(x) &= 2\ln\frac{\prod_i e^{-k_i} k_i^{k_i} \cdot k_i!}{\prod_ik_i! \cdot e^{-\hat{\mu_i}} \hat{\mu_i}^{k_i}} \\
	&= 2\ln\left[\prod_i e^{\hat{\mu_i}-k_i} \left(\frac{k_i}{\hat{\mu_i}}\right)^{k_i}\right] \\
	&= 2\sum_i\left[\hat{\mu_i}-k_i+k_i\ln\left(\frac{k_i}{\hat{\mu_i}}\right)\right] \Rightarrow \\
	\lambda(x) &= 2\sum_i\left[k_i\ln\left(\frac{k_i}{\hat{\mu_i}}\right)-\left(k_i-\hat{\mu_i}\right)\right]
\end{align*}

The above expression presents the GOF statistic for a generic histogram, which in the asymptotic limit is distributed like a $\chi^2_{DOF}$.

Before we move on, there is a rewarding connection to be made here. If you look closely at our final expression for $\lambda(x)$, it might not immediately look like the classic $\chi^2$ formula you likely encountered in your introductory courses, and that you certainly used to fit histograms in the past. However, if we assume our model is a reasonably good fit, such that the fluctuations $\delta_i = k_i - \hat{\mu}_i$ are relatively small, we can perform a Taylor expansion of the logarithm:
	
	\[
	\ln\left(\frac{k_i}{\hat{\mu}_i}\right) = \ln\left(1 + \frac{\delta_i}{\hat{\mu}_i}\right) \approx \frac{\delta_i}{\hat{\mu}_i} - \frac{\delta_i^2}{2\hat{\mu}_i^2} + \dots
	\]
	
Plugging this approximation back into our sum for $\lambda(x)$, and remembering that $\sum (k_i - \hat{\mu}_i) = \sum \delta_i$, a bit of algebra leads us to:
	
	\[
	\lambda(x) \approx \sum_i \frac{(k_i - \hat{\mu}_i)^2}{\hat{\mu}_i}
	\]
	
This is exactly the \textbf{Pearson's $\chi^2$ statistic}! It turns out that the LLR approach is effectively a more robust, fundamental version of the same idea. While the Pearson version is a second-order approximation, our LLR-based "Likelihood $\chi^2$" (often called the \emph{Deviance}) remains exact for any sample size, only relying on the asymptotic limit to claim its $\chi^2$ distribution. 

\subsection{Goodness of Fit of Normally distributed variables}

Though it can be perhaps seen as a particular case of the previous section, we take a moment to formally flush out the case of $N$ measurements $\{x_i\}_N$ of a Normally distributed variable with known variance $\sigma_i^2$ and unknown mean $\mu_i$. Both the mean and variance can be the same or different for each measure, though in this example the variance is always known. We want to verify the hypothesis $H_0$ that the means $\E\left[x_i\right]=\mu_i$. It's not hard to see how this is a specific case of the histogram problem for high-statistic bins, and so we start with a similar reasoning:

	\[
	\lambda(x) = 2\ln\frac{\sup_{\theta,\nu}\prod_i\N\left(x_i;\mu_i(\theta,\nu),\sigma_i\right)}{\sup_\theta\prod_i\N\left(x_i;\mu_i(\theta),\sigma_i\right)}
	\]

We have that $\sup_{\theta,\nu}\mu_i(\theta,\nu)=x_i$ itself, since the Likelihood is maximized when the expected value of the Normal distribution is exactly equal to its observed count. Therefore, we have:

\begin{align*}
	\lambda(x) &= 2\ln\frac{\prod_i\N\left(x_i;x_i,\sigma_i\right)}{\prod_i\N\left(x_i;\hat{\mu}_i,\sigma_i\right)} \\
	&= 2\ln\prod_i \frac{\sqrt{2\pi}\sigma_i \cdot e^{-\frac{(x_i-x_i)^2}{2\sigma_i^2}}}{\sqrt{2\pi}\sigma_i \cdot e^{-\frac{(x_i-\hat{\mu}_i)^2}{2\sigma_i^2}}} \Rightarrow \\
	\lambda(x) &= \sum_i \frac{\left(x_i-\hat{\mu}_i\right)^2}{\sigma_i^2}
\end{align*}

Look familiar? Once again we have found the \textbf{Pearson $\chi^2$ Test}, which we now see to be the point of convergence for histogram GOF testing in the case of "good" hypotheses, high statistics, or Normally distributed variables. We can also see that in this case, the statistic has \emph{exactly} the distribution of a $\chi^2_{DOF}$.

\subsubsection{Unknown variance Normally distributed GOF}
Let's follow the logic of the previous point, but introduce the complication of an unknown variance. We are testing whether the data is consistent with a Normal distribution of known mean $\mu_0$ against a general Normal distribution where both mean and variance are free. To state this clearly, we have $H_0: p(x) = \mathcal{N}(x;\mu_0, \sigma)$ against $H_1: p(x) = \mathcal{N}(x;\mu, \sigma)$. This case is remarkably common in instrumental calibration, where the "true" value $\mu_0$ is known from a reference standard, but the precision of the instrument ($\sigma$) is unknown and must be estimated from the data itself.

We therefore introduce the unrestricted and restricted maximum likelihood estimators of the variance:

	\[
	\hat{\sigma}^2 = \frac{\sum_j\left(x_j-\bar{x}\right)^2}{N}\text{ and } \hat{\sigma}_0^2 = \frac{\sum_j\left(x_j-\mu_0\right)^2}{N}
	\]

Applying our trusty LLR:

\begin{align*}
	\lambda(x) &= 2\ln \frac{\sup_{\mu,\sigma}\prod_i^N\N\left(x_i,\mu, \sigma\right)}{\sup_{\sigma}\prod_i^N\N\left(x_i,\mu_0, \sigma\right)} \\ 
	&= 2\ln \frac{\prod_i^N\N\left(x_i,\bar{x}, \hat\sigma\right)}{\prod_i^N\N\left(x_i,\mu_0, \hat{\sigma}_0\right)} \\ 
	&= 2\ln \frac{\prod_i^N\sqrt{2\pi}\hat{\sigma}_0 \cdot e^{-\frac{\left(x_i-\bar{x}\right)^2}{2\hat{\sigma}^2}}}{\prod_i^N\sqrt{2\pi}\hat{\sigma} \cdot e^{-\frac{\left(x_i-\bar{x}\right)^2}{2\hat{\sigma}_0^2}}} \\ 
	&= 2\ln\prod_i^N\left(\frac{\hat{\sigma}_0}{\hat{\sigma}}\right) + \sum_i^N\frac{\left(x_i-\mu_0\right)^2}{2\hat{\sigma}_0^2}-\sum_i^N\frac{\left(x_i-\bar{x}\right)^2}{2\hat{\sigma}^2}
\end{align*}

Looking at our above definition of the estimators for the variance, we can invert them to write 

	\[
	\sum_i^N\left(x_i-\bar{x}\right)^2=N\hat{\sigma}^2\text{ and }\sum_i^N\left(x_i-\mu_0\right)^2=N\hat{\sigma}_0^2
	\]

Therefore, the two sum terms in our expression both equal $\frac{N}{2}$ and cancel out, leaving us with:

	\[
	\lambda(x) = 2\ln\prod_i^N\left(\frac{\hat{\sigma}_0}{\hat{\sigma}}\right) = N\ln\left(\frac{\hat{\sigma}_0^2}{\hat{\sigma}^2}\right)
	\]

Using that the ratio is constant across observations, we have:
	\[
	\lambda(x) = N \ln\frac{1/N\sum_j\left(x_j-\mu_0\right)^2}{1/N\sum_j\left(x_j-\bar{x}\right)^2}\\
	\]

We expand the numerator term:

	\[
	\sum_j\left(x_j-\mu_0\right)^2 = \sum_j\left(x_j+\bar{x}-\bar{x}-\mu_0\right)^2
	\]

Which we can re-write in the following way with a little algebra:

	\[
	\sum_j\left(x_j-\bar{x}\right)^2 + \sum_j\left(\bar{x}-\mu_0\right)^2 + 2\sum_j\left(x_j-\bar{x}\right)\left(\bar{x}-\mu_0\right)
	\]

The final mixed term is equal to zero: $\sum_j\left(x_j-\bar{x}\right)=0$, and we are left with:

	\[
	\sum_j\left(x_j-\mu_0\right)^2 = \sum_j\left(x_j-\bar{x}\right)^2 + N\left(\bar{x}-\mu_0\right)^2
	\]

Substituting back into the LLR, we have:

\begin{align*}
	\lambda(x) &= N \ln \left[\frac{\sum (x_i - \bar{x})^2 + N(\bar{x} - \mu_0)^2}{\sum (x_i - \bar{x})^2} \right] \\
	&= N \ln \left[1 + \frac{N(\bar{x} - \mu_0)^2}{\sum (x_i - \bar{x})^2} \right] \\
\end{align*}

We'd like to use our old friend Taylor to approximate the logarithm, but we must first convince ourselves that the second term is much smaller than $1$, at least asymptotically.
First of all, let's identify that the expected value of $\left(\bar{x}-\mu_0\right)^2$ is, by definition under $H_0$, the variance of $\bar{x}$, which is $\frac{\sigma^2}{N}$. This means that $\left(\bar{x}-\mu_0\right)^2$ is typically of order $\frac{1}{N}$ ($\mathcal{O}_p(\frac{1}{N})$ \emph{in probability}).

Looking at the rest of the term, we can identify $\frac{N}{\sum (x_i - \bar{x})^2} = \frac{1}{\hat{\sigma}^2}$, the estimator for the variance, which scales as $\mathcal{O}(1)$. Putting this together, we can see that:

	\[
	\frac{N(\bar{x} - \mu_0)^2}{\sum (x_i - \bar{x})^2} \sim \mathcal{O}\left(N^{-1}\right)
	\]

Using the asymptotic expansion $\ln(1+x) \approx x$, we see that:

	\[
	\lambda(x) \approx N^2\frac{\left(\bar{x}-\mu_0\right)^2}{\sum_j\left(x_j-\bar{x}\right)^2}
	\]

Let's call $s$ (the \emph{unbiased} sample variance) $s^2 = \frac{\sum_i\left(x_i-\bar{x}\right)^2}{N-1}$. We then have that

	\[
	\lambda(x) \approx \frac{N}{N-1}\underbrace{\frac{N\left(\bar{x}-\mu_0\right)^2}{s^2}}_{t^2}
	\]

Where we have defined $t$ to be a specific Test that we will talk about in the next section. In this more complex case, the likelihood-ratio test is equivalent to the $t$-test. For large $N$, the factor $\frac{N}{N-1} \to 1$, and $\lambda(x)\approx t^2$, which we will learn to converge in distribution to $\chi^2_1$.

\section{Student's t}
Annoying derivation, I already did the one above I can't handle two in the same day
\section{Continuous distributions}

\section{1-d}

Of course all of this has been randomly sectioned, we can touch up on that as we write. 